% !TEX program = pdflatex
% The AAAI 2027 style requires pdfTeX; do not compile this paper with XeLaTeX.
\documentclass[letterpaper]{article} % DO NOT CHANGE THIS
\usepackage[submission]{aaai2027}  % DO NOT CHANGE THIS
\usepackage[hyphens]{url}  % DO NOT CHANGE THIS
\usepackage{graphicx} % DO NOT CHANGE THIS
\urlstyle{rm} % DO NOT CHANGE THIS
\def\UrlFont{\rm}  % DO NOT CHANGE THIS
\usepackage{natbib}  % DO NOT CHANGE THIS AND DO NOT ADD ANY OPTIONS TO IT
\usepackage{caption} % DO NOT CHANGE THIS AND DO NOT ADD ANY OPTIONS TO IT
\frenchspacing  % DO NOT CHANGE THIS

\usepackage{algorithm}
\usepackage{algorithmic}
\usepackage{booktabs}
\usepackage{multirow}
\usepackage{amsmath,amssymb}

\pdfinfo{
/TemplateVersion (2027.1)
}

\setcounter{secnumdepth}{2}

\title{Overlap-Aware Selective Forgetting in Continual Learning\\via Protected Subnetworks}
\author{
    Anonymous Submission
}
\affiliations{
}

\begin{document}

\maketitle

\begin{abstract}
Continual learning systems deployed in real-world environments must often honor requests to unlearn specific information. In task-incremental settings, selective forgetting requires suppressing a designated target task while maintaining performance on all retained tasks. However, when tasks share underlying model parameters, resetting or repairing weights for one task frequently degrades others. To address this parameter overlap, we introduce EPALL, an enhanced Privacy-Aware Lifelong Learning approach. Our approach identifies shared parameter coordinates, ranks them according to their importance to retained tasks, and anchors high-sensitivity weights during the repair phase. Evaluated on Split-CIFAR benchmarks, EPALL effectively reduces forgotten-task accuracy to the chance level while staying within 0.5 percentage points of full-network baselines. Compared to the original PALL algorithm, it consistently achieves higher final accuracy and significantly reduces the maximum drop in task performance. Additionally, we analyze an alternative parameter-efficient adapter mechanism showing that unlearning in this setup is primarily driven by target-path resets. Overall, EPALL provides a robust framework for empirical selective forgetting in shared-parameter continual learning.
\end{abstract}

% Insert the verified anonymous code-archive URL before submission; keep blank until final packaging.
% \begin{links}
%     \link{Code}{https://anonymous.4open.science/...}
% \end{links}

\section{Introduction}

Continual learning (CL) systems acquire tasks sequentially while trying to prevent catastrophic forgetting \cite{mccloskey1989catastrophic,kirkpatrick2017overcoming}. Selective forgetting poses the inverse operational problem: after a task-level deletion request, the model should suppress that task without catastrophically degrading the tasks that remain active \cite{cao2015towards,bourtoule2021machine}. We study this task-incremental setting and do not equate task suppression with deleting individual records, raw data, checkpoints, or external copies.

The core difficulty is \emph{parameter overlap}: a high-capacity network spreads each task's knowledge across many shared weights. Running an unconstrained unlearning step on a target task $D_f$---gradient ascent, or Fisher-based scrubbing \cite{golatkar2020forgetting}---therefore also moves weights that the retained tasks $D_r$ depend on, and their accuracy drops as a side effect.

The Privacy-Aware Lifelong Learning (PALL) framework \cite{ozdenizci2025pall} modularizes task knowledge with stored subnetworks, but its request-time reset and repair do not distinguish shared coordinates by their sensitivity to retained tasks---which matters under heavy overlap, where perturbing a small but important shared region damages several active tasks.

\paragraph{Contributions.}
\begin{enumerate}
    \item \textbf{Retention-sensitive shared-weight protection:} EPALL measures structural overlap between the forgotten and active subnetworks, ranks shared coordinates by retained-task sensitivity, and anchors the critical subset during repair.
    \item \textbf{Matched evaluation and a controlled overlap benchmark:} an eleven-method, three-seed Split-CIFAR comparison reports retention, forgetting, target accuracy, and request time jointly with paired analysis; a further $120$-run benchmark sweeps five pre-specified overlap levels across six methods.
    \item \textbf{Negative mechanism audit:} controls over a parameter-efficient adapter pathway isolate target reset, repair, soft masking, and classifier ascent, and attribute its suppression to reset rather than the overlap-aware update---so a reset-driven effect is not presented as a second contribution.
\end{enumerate}

\section{Related Work}

\subsection{Continual Learning and Memory Replay}
Continual learning accumulates knowledge across a stream of tasks \cite{thrun1998lifelong,delange2022survey}. Regularization methods anchor parameters \cite{kirkpatrick2017overcoming,zenke2017continual,li2017learning}, while rehearsal methods keep an episodic buffer \cite{buzzega2020dark,rebuffi2017icarl}. The gradient-projection variants GEM and A-GEM \cite{lopezpaz2017gem,chaudhry2019agem} are closest in spirit to our mechanism: they also let retained-task gradients constrain the update, but project it globally rather than protecting a selected shared region. A third family isolates parameters per task via pruning \cite{mallya2018packnet}, hard attention masks \cite{serra2018hat}, or supermasks over a fixed backbone \cite{wortsman2020supsup}, building on the fact that trainable subnetworks exist inside dense or even random networks \cite{frankle2019lottery,ramanujan2020hidden}. PALL inherits its per-task binary masks from this line, specifically winning-subnetwork continual learning \cite{kang2022wsn}, which is what makes \emph{structural} overlap between task masks well defined. CLPU \cite{liu2022clpu} is the strongest isolation-based reference: each active task gets its own full network and a request deletes that network, so interference is impossible by construction but storage grows with the task count. EPALL targets the complementary case, where tasks already share one model.

\subsection{Machine Unlearning Strategies}
Retraining-based methods partition data into isolated ensemble members \cite{bourtoule2021machine}, and amnesiac training stores per-batch updates for later subtraction \cite{graves2021amnesiac}. A parallel theoretical line gives deletion guarantees for convex models \cite{guo2020certified,neel2021descent,sekhari2021remember}; these do not transfer to the non-convex sequential setting we study, which is why our claims stay empirical. Approximate methods instead edit gradients to scrub information \cite{golatkar2020forgetting}, and Boundary Unlearning shifts output decision boundaries for speed \cite{chen2023boundary}. In a continual setting, though, an unconstrained update can damage feature space that active tasks still use. Saliency Unlearning (SalUn) \cite{fan2024salun} confines updates to influential weights; Selective Synaptic Dampening (SSD) \cite{foster2024ssd} dampens parameters by Fisher-based importance. Controlled distillation can jointly support learning and removal with a fixed memory \cite{chatterjee2024unified}. UnCLe instead generates task-specific parameters and explicitly studies post-deletion task relapse without retained data \cite{adhikari2026uncle}; CATA studies sequential removal and persistence for vision--language models through conflict-averse task-vector aggregation \cite{lin2026cata}. Our task-incremental shared-model setting keeps a bounded rehearsal memory and edits overlapping weights directly, but these works motivate a separate persistence audit rather than treating immediate post-request accuracy as sufficient.

\subsection{Parameter-Efficient Fine-Tuning}
Parameter-efficient fine-tuning (PEFT), through adapters \cite{houlsby2019parameter} and LoRA \cite{hu2021lora}, restricts training to compact modules inserted into a frozen network. The recent BID-LoRA preprint \cite{rachapudi2026bidlora} explores a related continual-learning/unlearning combination. We audit whether measuring overlap inside a shared bottleneck changes request-time endpoints beyond target-path reset and retained-task repair.

\section{Methodology}

\subsection{The Limitation of Unconstrained Forgetting}
PALL identifies task-specific parameters, resets or reconstructs coordinates from stored masks, and repairs the retained tasks. What it does not do is separate shared coordinates by how sensitive the retained tasks are to them. Under heavy overlap, reset and repair can therefore disturb pathways that several active tasks lean on. We map that intersection and protect it directly.

\subsection{EPALL: Retention-Sensitive Shared-Weight Protection}
Let $M_f$ be the binary subnetwork mask stored for the task requested for deletion and let $M_r$ be the union of the masks of the retained tasks. Their shared support is
\begin{equation}
S_{share}=M_f\cap M_r.
\end{equation}
EPALL ranks $S_{share}$ by retained-buffer gradient magnitude (or gradient-conflict energy in an ablation) and selects $S_{share\_crit}$. These binary-subnetwork sets differ from the gradient-selected adapter sets below. The target rehearsal entry is used only during request-time diagnostics and is then removed from in-process memory; no raw-data, checkpoint, or external-copy deletion is claimed. The inherited PALL reset is branch-dependent because retained backups and historical masks determine which coordinates are preserved or reconstructed. Retention repair is restricted to the nominal shared repair mask and minimizes
\begin{equation}
\label{eq:repair}
\begin{aligned}
\mathcal{L}_{repair} ={}& \alpha\,\mathcal{L}_{DER} + \beta\,\mathcal{L}_{CE}\\
&+ \lambda \sum_{w \in S_{share\_crit}} (w-w_{anchor})^2 .
\end{aligned}
\end{equation}
Here $\mathcal{L}_{DER}$ matches stored logits, $\mathcal{L}_{CE}$ uses retained labels, and the final term limits drift on the selected critical shared coordinates. The default anchor is the pre-request value $w_{old}$; we separately evaluate a fresh reinitialization anchor $w_{reinit}$ because $w_{old}$ can itself encode information about the forgotten task.

\begin{algorithm}[t]
\caption{EPALL Selective Forgetting}
\label{alg:pall-modified}
\textbf{Input}: target task $f$, task masks $\{M_t\}$, rehearsal memory $\mathcal{B}$, model weights $\theta$, protection ratio $\rho$, anchor weight $\lambda$
\begin{algorithmic}[1]
\STATE Copy $M_f$; if conflict scoring is selected, cache its forget gradient while $\mathcal{B}_f$ is available
\STATE Delete $\mathcal{B}_f$ and remove $M_f$ from the active mask set
\STATE $M_r \leftarrow \bigvee_{t\ne f}M_t$; $S_{share}\leftarrow M_f\cap M_r$
\STATE Compute retained-buffer sensitivity $q_i=|\nabla_i\mathcal{L}_{retain}|$ for $i\in S_{share}$
\STATE $S_{share\_crit}\leftarrow$ the top $\lceil\rho|S_{share}|\rceil$ coordinates by $q_i$
\STATE Cache anchors on $S_{share\_crit}$ and sparse retained-task backups required by the inherited reset branches
\STATE Apply PALL's branch-dependent target reset/reconstruction; derive the nominal retained-repair mask
\FOR{$k=1,\ldots,N_{repair}$}
  \STATE Sample retained rehearsal entries and minimize Eq.~\eqref{eq:repair} only on the repair mask
\ENDFOR
\STATE Rebuild the active union mask and discard request-only target state
\end{algorithmic}
\end{algorithm}

Algorithm~\ref{alg:pall-modified} states the default gradient-magnitude path used in the main experiments, and Figure~\ref{fig:epall} summarizes the problem, the selection rule, and the request-time procedure. The conflict ablation changes only line 4, while the direct component study changes the anchor or the candidate/ranking rule without changing the request schedule.

\begin{figure*}[t]
\centering
\includegraphics[width=.92\textwidth]{Figures/EPALL_mechanism.png}
\caption{EPALL. (a) The problem: forgotten- and retained-task paths run through shared coordinates (orange), so a naive reset causes collateral retained-task drift. (b) Selection: the candidate set is the structural overlap $S_{share}=M_f\cap M_r$ between the forgotten mask and the retained-mask union; within it the top-$\rho$ coordinates ranked by retained-buffer gradient sensitivity form the protected subset $S_{share\_crit}$. (c) The request, in implementation order: the target buffer $\mathcal{B}_f$ is deleted at the request boundary once its optional diagnostic has been computed; the retained union is then formed, the overlap ranked, and anchors cached on $S_{share\_crit}$; PALL's branch-dependent reset suppresses the target path; and retained-only repair under Eq.~\eqref{eq:repair} holds the protected subset near its cached anchor. The objectives are empirical, not certified-removal guarantees.}
\label{fig:epall}
\end{figure*}

\subsection{Exploratory Parameter-Efficient Adapter Pathway}
% Preserve the advisor's prose edit; do not convert these three groups back to a list.
To cut the update footprint and keep tasks from interfering, we freeze the feature-extraction layers ($\theta_{base}$). Only the following three groups stay trainable: shared bottleneck adapter ($\phi_{shared}$), task-specific adapters ($\phi_{t}$), and one shared classifier, used through task-specific row slices. Both adapters are standard post-backbone residual bottlenecks with ReLU \cite{houlsby2019parameter,nair2010relu}, applied in sequence (shared, then task-specific).

A request follows a strict order so that every forget-dependent gradient is estimated while the target rehearsal entry is still available. First, absolute shared-adapter gradients on the target and retained buffers define top-ratio sets $S_{forget}$ and $S_{active}$; their intersection is $S_{share\_crit}$, and its top $\lceil\alpha_p|S_{share\_crit}|\rceil$ coordinates by retained-gradient magnitude (gradient-conflict energy only in the stated ablation) form the hard-protected set $H$. Second, the target adapter is reinitialized (Kaiming down-projection, zero up-projection), its classifier slice is reset, and the shared adapter is optimized toward a uniform distribution over the target task's classes with each coordinate's step scaled by
\begin{equation}
\label{eq:softmask}
m_{soft}(w) = \begin{cases}
0, & w \in H, \\
1-p, & w \in S_{share\_crit}\setminus H, \\
1, & w \in S_{forget}\setminus S_{active}, \\
0, & \text{otherwise,}
\end{cases}
\end{equation}
with $p\in[0,1]$. Third, the target classifier receives a gradient-ascent update from its pre-reset forget gradient, the target rehearsal entry is deleted, and all trainable parameters are briefly repaired on retained buffers.

The mask therefore governs only the shared-adapter displacement during the uniform-target subphase; it is not a permanent protection mask, and it does not bound end-to-end retained accuracy, because the reset, the classifier ascent, and the unmasked repair also move the endpoint. The component controls below test it empirically rather than treating the multiplier as a guarantee. Since the final shared adapter remains a function of target-buffer gradients, this pathway is neither exact-unlearning nor privacy-certified.

% Declared before the setup text so the full-width primary table can occupy the
% next available page top without leaving a half-empty column.
% AUTO-GENERATED by tools/generate_main_standard_table.py; do not edit.
% Source: server_report_table.csv
\begin{table*}[t]
\centering
\footnotesize
\setlength{\tabcolsep}{2.2pt}
\begin{tabular}{@{}llccccc@{}}
\toprule
\multicolumn{7}{@{}l}{\textit{Panel A: Standard Split-CIFAR-10}} \\
\cmidrule(lr){1-7}
\textbf{Family} & \textbf{Method} & \textbf{$A_{\mathrm{final}}\!\uparrow$} & \textbf{$F_{\mathrm{avg}}\!\downarrow$} & \textbf{WorstDrop $\downarrow$} & \textbf{$A_u\!\rightarrow\!$ chance} & \textbf{$T_f\!\downarrow$ (s)} \\
\midrule
\multirow{1}{*}{\textit{Isolation}} & CLPU & $0.9604\mathbin{\pm}0.0191$ & $0.2949\mathbin{\pm}0.0414$ & $0.0000\mathbin{\pm}0.0000$ & $0.5210\mathbin{\pm}0.0048$ & 0.0005 \\
\addlinespace[2pt]
\multirow{4}{*}{\textit{Shared CL}} & DER++ & $0.9191\mathbin{\pm}0.0356$ & $0.2702\mathbin{\pm}0.0804$ & $0.0037\mathbin{\pm}0.0032$ & $0.5225\mathbin{\pm}0.1032$ & 2.3188 \\
 & ER & $0.9054\mathbin{\pm}0.0399$ & $0.3061\mathbin{\pm}0.0464$ & $0.0123\mathbin{\pm}0.0168$ & $0.4485\mathbin{\pm}0.0495$ & 2.0622 \\
 & EWC & $0.8436\mathbin{\pm}0.0724$ & $0.2964\mathbin{\pm}0.1496$ & $0.0000\mathbin{\pm}0.0000$ & $0.9418\mathbin{\pm}0.0298$ & 0.0001 \\
 & LwF & $0.9406\mathbin{\pm}0.0185$ & $0.0866\mathbin{\pm}0.0604$ & $0.0000\mathbin{\pm}0.0000$ & $0.9582\mathbin{\pm}0.0196$ & 0.0000 \\
\addlinespace[2pt]
\multirow{2}{*}{\textit{Unlearning}} & SSD & $0.9011\mathbin{\pm}0.0581$ & $0.2241\mathbin{\pm}0.1560$ & $0.0027\mathbin{\pm}0.0025$ & $0.8030\mathbin{\pm}0.2715$ & 0.2330 \\
 & SalUn & $0.9032\mathbin{\pm}0.0114$ & $0.3077\mathbin{\pm}0.0743$ & $0.0100\mathbin{\pm}0.0087$ & $0.5452\mathbin{\pm}0.0294$ & 2.1656 \\
\addlinespace[2pt]
\multirow{2}{*}{\textit{Full PALL}} & PALL-Original & $0.9513\mathbin{\pm}0.0118$ & $0.2984\mathbin{\pm}0.0339$ & $0.0067\mathbin{\pm}0.0060$ & $0.5000\mathbin{\pm}0.0000$ & 2.1345 \\
 & \textbf{EPALL} & $0.9555\mathbin{\pm}0.0149$ & $0.2956\mathbin{\pm}0.0362$ & $0.0035\mathbin{\pm}0.0044$ & $0.5000\mathbin{\pm}0.0000$ & 2.7933 \\
\addlinespace[2pt]
\multirow{2}{*}{\textit{PEFT$^\dagger$}} & LoRA$^\dagger$ & $0.7863\mathbin{\pm}0.0810$ & $0.1687\mathbin{\pm}0.0442$ & $0.0035\mathbin{\pm}0.0031$ & $0.5000\mathbin{\pm}0.0005$ & 3.2665 \\
 & \textbf{PALL-Adapter}$^\dagger$ & $0.7822\mathbin{\pm}0.0753$ & $0.1695\mathbin{\pm}0.0341$ & $0.0075\mathbin{\pm}0.0067$ & $0.4927\mathbin{\pm}0.0127$ & 3.5084 \\
\addlinespace[4pt]
\midrule
\addlinespace[2pt]
\multicolumn{7}{@{}l}{\textit{Panel B: Standard Split-CIFAR-100}} \\
\cmidrule(lr){1-7}
\textbf{Family} & \textbf{Method} & \textbf{$A_{\mathrm{final}}\!\uparrow$} & \textbf{$F_{\mathrm{avg}}\!\downarrow$} & \textbf{WorstDrop $\downarrow$} & \textbf{$A_u\!\rightarrow\!$ chance} & \textbf{$T_f\!\downarrow$ (s)} \\
\midrule
\multirow{1}{*}{\textit{Isolation}} & CLPU & $0.7391\mathbin{\pm}0.0079$ & $0.1942\mathbin{\pm}0.0165$ & $0.0000\mathbin{\pm}0.0000$ & $0.1020\mathbin{\pm}0.0044$ & 0.0007 \\
\addlinespace[2pt]
\multirow{4}{*}{\textit{Shared CL}} & DER++ & $0.6718\mathbin{\pm}0.0271$ & $0.2570\mathbin{\pm}0.0057$ & $0.0357\mathbin{\pm}0.0081$ & $0.2297\mathbin{\pm}0.0788$ & 10.7241 \\
 & ER & $0.6059\mathbin{\pm}0.0407$ & $0.2726\mathbin{\pm}0.0232$ & $0.0740\mathbin{\pm}0.0401$ & $0.3553\mathbin{\pm}0.0538$ & 6.2048 \\
 & EWC & $0.4685\mathbin{\pm}0.0424$ & $0.4002\mathbin{\pm}0.0220$ & $0.0000\mathbin{\pm}0.0000$ & $0.4487\mathbin{\pm}0.0975$ & 0.0001 \\
 & LwF & $0.6938\mathbin{\pm}0.0390$ & $0.1440\mathbin{\pm}0.0105$ & $0.0000\mathbin{\pm}0.0000$ & $0.7743\mathbin{\pm}0.0790$ & 0.0000 \\
\addlinespace[2pt]
\multirow{2}{*}{\textit{Unlearning}} & SSD & $0.6381\mathbin{\pm}0.0212$ & $0.3048\mathbin{\pm}0.0385$ & $0.0040\mathbin{\pm}0.0046$ & $0.2977\mathbin{\pm}0.0870$ & 0.4086 \\
 & SalUn & $0.6191\mathbin{\pm}0.0449$ & $0.2681\mathbin{\pm}0.0262$ & $0.0750\mathbin{\pm}0.0325$ & $0.3333\mathbin{\pm}0.0241$ & 6.1575 \\
\addlinespace[2pt]
\multirow{2}{*}{\textit{Full PALL}} & PALL-Original & $0.7155\mathbin{\pm}0.0155$ & $0.2116\mathbin{\pm}0.0139$ & $0.0210\mathbin{\pm}0.0040$ & $0.1037\mathbin{\pm}0.0055$ & 10.8251 \\
 & \textbf{EPALL} & $0.7346\mathbin{\pm}0.0095$ & $0.1989\mathbin{\pm}0.0146$ & $0.0010\mathbin{\pm}0.0026$ & $0.0903\mathbin{\pm}0.0085$ & 14.8144 \\
\addlinespace[2pt]
\multirow{2}{*}{\textit{PEFT$^\dagger$}} & LoRA$^\dagger$ & $0.2686\mathbin{\pm}0.0056$ & $0.0595\mathbin{\pm}0.0139$ & $0.0170\mathbin{\pm}0.0170$ & $0.1097\mathbin{\pm}0.0085$ & 7.0411 \\
 & \textbf{PALL-Adapter}$^\dagger$ & $0.2643\mathbin{\pm}0.0255$ & $0.0772\mathbin{\pm}0.0101$ & $0.0290\mathbin{\pm}0.0165$ & $0.1027\mathbin{\pm}0.0031$ & 8.1277 \\
\bottomrule
\end{tabular}
\caption{Matched standard Split-CIFAR comparison (20 epochs/task; three seeds). Rows are grouped by operating assumption; bold method names denote our methods, while numeric cells are not ranked because the metrics must be read jointly. CLPU is the full-model isolation reference. $^\dagger$The PEFT rows freeze a random backbone here; their intended pretrained setting appears in Table~\ref{tab:pretrained}. Chance $A_u$ is $0.5$ for CIFAR-10 and $0.1$ for CIFAR-100. All metrics except the non-hardware-normalized $T_f$ are mean${\pm}$std.}
\label{tab:main_standard}
\end{table*}

\section{Experimental Setup}

We evaluate task-level continual forgetting over three public benchmarks with fixed schedules that interleave task-training and forgetting requests. The benchmarks span increasing task count and representation difficulty: CIFAR-10 \cite{krizhevsky2009learning} as 5 two-class tasks (chance $0.5$); CIFAR-100 \cite{krizhevsky2009learning} as 10 five-class tasks in the overlap-heavy protocol (chance $0.2$) and 10-class tasks in the standard protocol (chance $0.1$); and TinyImageNet \cite{le2015tinyimagenet} as 20 ten-class tasks (chance $0.1$).
Apart from the ViT ablation, CIFAR-10, TinyImageNet, and pretrained-frozen runs use ResNet-18 variants; main CIFAR-100 runs use ResNet-34 variants \cite{he2016deep}. Configurations use SGD \cite{robbins1951sgd} (learning rate $0.01$, momentum $0.9$, weight decay $5\times10^{-4}$, batch size $32$) and a replay budget of $500$, partitioned a priori across tasks without post-deletion reallocation. Split-CIFAR-100 LoRA alone uses $10^{-3}$ because $10^{-2}$ diverges. Main schedules train each task at most once, contain three forget requests, and never reintroduce a deleted task; the standard, overlap-heavy, and pretrained protocols use 20 epochs, 3 epochs, and a frozen ImageNet backbone respectively. PALL-Adapter uses task/shared bottlenecks $16$, $\alpha_f{=}0.3$, $\alpha_p{=}0.2$, $p=|S_{share\_crit}|/|S_{forget}|$, 10 forget steps, and 50 repair steps. EPALL uses protection ratio $0.2$ and $\lambda{=}1$; LoRA rank is 8; subnet sparsity is $0.8$ on CIFAR-10/TinyImageNet and $0.9$ on CIFAR-100. Reruns used one NVIDIA RTX 3070 Ti (8\,GB), dual Xeon E5-2620 v4 CPUs, 31\,GB RAM, Ubuntu 24.04.4, Python 3.12.3, PyTorch 2.5.1, and deterministic scripts. The standard, pretrained-PEFT, overlap-heavy, and component-control tables use seeds $\{0,1,2\}$; the anchor and persistence analyses use $\{0,1\}$. These counts support descriptive, not confirmed, rankings.

We report retained-task $A_{\mathrm{final}}$ ($\uparrow$). For CL context, $F_{\mathrm{avg}}(r)=|\mathcal{T}_{trained}|^{-1}\sum_t[\max_{k<r}A_t(k)-A_t(r)]$ covers every task trained by request $r$, including later deletions; request-specific claims use retained-only Fu and WorstDrop plus forgotten-task Au. Fu is retained mean before minus after a request, and
\begin{equation}
\label{eq:worstdrop}
\mathrm{WorstDrop} = \max_{t \in \mathcal{T}_{active}} \left( A_{t}^{before} - A_{t}^{after} \right)
\end{equation}
($\downarrow$) is signed, not clipped. Au is measured immediately and read through $|\mathrm{Au}-\mathrm{chance}|$. Tables use each run's final forget event, then mean and sample standard deviation across seeds. $R_{\mathrm{upd}}$ is only the nominal retained-repair trainable fraction, not total compute, storage, memory, runtime, or changed coordinates.

\section{Experimental Results and Analysis}

\subsection{Matched Standard Split-CIFAR Comparison}
\label{subsec:standard-protocol}
Table~\ref{tab:main_standard} is the primary comparison. It uses 20 epochs per task and three seeds; within each seed, all eleven methods receive the same standard Split-CIFAR request schedule. SSD and SalUn come from their matched-protocol aggregates; the separate single-seed SSD tuning runs are not used.

CLPU is the strongest row overall and we do not claim to outperform it: the evaluated implementation stores one full network per active task and forgets by deleting that entry, so its sub-millisecond $T_f$ times a dictionary deletion rather than the cost of maintaining those networks. It is attractive when such isolation is acceptable, but it is not a shared-model solution.

Among the shared-model methods, EPALL holds $A_{\mathrm{final}}$ within half a point of CLPU ($0.9555$ vs.\ $0.9604$ on CIFAR-10; $0.7346$ vs.\ $0.7391$ on CIFAR-100), puts $A_u$ at or near chance ($0.5000$ and $0.0903$), and keeps WorstDrop at $0.0035$ and $0.0010$. It thus approaches the isolation reference without giving every task its own network. The other columns matter too: EWC and LwF post zero WorstDrop on CIFAR-10 but leave $A_u$ at $0.9418$ and $0.9582$, preserving the task rather than forgetting it, while SSD leaves $A_u=0.8030$ on CIFAR-10 and both SSD and SalUn stay well above the $0.1$ target on CIFAR-100.

Since schedules are paired by seed we also compare EPALL with PALL-Original directly. On CIFAR-100 $\Delta A_{\mathrm{final}}$ is positive in all three pairs (mean $+0.0191$, bootstrap $[0.0103,0.0290]$), with mean gains of $0.0127$ in $F_{\mathrm{avg}}$ and $0.0200$ in WorstDrop; on CIFAR-10 the accuracy gain is smaller ($+0.0043$, $[0.0000,0.0065]$). Retention costs request time: EPALL is ${\approx}4.0$\,s slower than PALL-Original on CIFAR-100.

Exact one-sided paired tests expose the limits of a three-seed design: the smallest attainable $p$ for both the Wilcoxon signed-rank and sign tests is $2^{-3}=0.125$, so no comparison here can reach a conventional threshold. CIFAR-100 hits exactly that floor on all three metrics---every pair favours EPALL, $d_z=2.0$, $12.0$, $7.6$---while CIFAR-10 gives $p=0.25$--$0.5$ and $d_z\le1.2$ because pairs tie. We therefore claim direction consistency with large paired effects, not significance.

The two daggered PEFT rows are diagnostics, not capacity-matched competitors: under this from-scratch protocol LoRA and PALL-Adapter both sit on a randomly initialized frozen ResNet. Their CIFAR-100 accuracies are nearly identical ($0.2686$ and $0.2643$) despite different forgetting mechanisms, and both recover once the backbone is pretrained (Table~\ref{tab:pretrained}), which points to a representation bottleneck rather than the forgetting step.

\subsection{EPALL Mechanism Attribution}
\label{subsec:modified-components}
Table~\ref{tab:modified-components} asks whether the reported gains actually come from the proposed mechanism. Five matched arms share the overlap-heavy protocol, schedules, and seeds and differ in one ingredient each: \emph{no-anchor} sets $\lambda=0$; \emph{overlap-only} anchors the entire structural overlap without ranking; \emph{random-budget} anchors randomly chosen structural-overlap coordinates with the same count as full; \emph{ranking-no-overlap} ranks by retained gradients without the explicit overlap filter; \emph{full} is the deployed method.

\begin{table}[t]
\centering
\footnotesize
\setlength{\tabcolsep}{1.4pt}
\begin{tabular}{@{}llccc@{}}
\toprule
\textbf{Data} & \textbf{Control} & \textbf{$A_{\mathrm{final}}$} & \textbf{W.Drop} & \textbf{$|A_u-c|$} \\ \midrule
\multirow{5}{*}{C10}
 & No anchor & .9128${\pm}$.0273 & .0033${\pm}$.0045 & .0000${\pm}$.0000 \\
 & Overlap only & .9145${\pm}$.0288 & $-$.0002${\pm}$.0003 & .0000${\pm}$.0000 \\
 & Random budget & .9111${\pm}$.0256 & .0050${\pm}$.0043 & .0000${\pm}$.0000 \\
 & Ranking/no-ovl & .9140${\pm}$.0281 & .0015${\pm}$.0015 & .0000${\pm}$.0000 \\
 & Full & .9140${\pm}$.0281 & .0015${\pm}$.0015 & .0000${\pm}$.0000 \\ \addlinespace[2pt]
\multirow{5}{*}{C100}
 & No anchor & .4965${\pm}$.0083 & .0093${\pm}$.0095 & .0087${\pm}$.0133 \\
 & Overlap only & .4917${\pm}$.0102 & .0013${\pm}$.0012 & .0027${\pm}$.0046 \\
 & Random budget & .4995${\pm}$.0097 & .0287${\pm}$.0083 & .0000${\pm}$.0000 \\
 & Ranking/no-ovl & .5119${\pm}$.0069 & .0093${\pm}$.0050 & .0007${\pm}$.0012 \\
 & Full & .5119${\pm}$.0069 & .0093${\pm}$.0050 & .0007${\pm}$.0012 \\ \bottomrule
\end{tabular}
\caption{Direct three-seed EPALL mechanism controls under the overlap-heavy protocol with matched schedules. Each arm changes one ingredient of the protection mechanism; no cells are ranked.}
\label{tab:modified-components}
\end{table}

On CIFAR-100 the full mechanism is the operative ingredient rather than a bystander: removing the anchor lowers paired $A_{\mathrm{final}}$ by $0.0154$ (descriptive bootstrap interval $[0.0057,0.0277]$), anchoring the whole overlap without ranking lowers it by $0.0202$ $[0.0066,0.0317]$, and spending the same anchor budget on random overlap coordinates lowers accuracy by $0.0124$ $[0.0097,0.0163]$ while raising WorstDrop by $0.0193$ $[0.0080,0.0340]$; all three intervals exclude zero and the signs agree in all three schedule pairs. Two further readings follow. First, a positive finding about \emph{where} the value lies: ranking-no-overlap reproduces full exactly, so the explicit intersection is not what carries the effect---retained-gradient ranking already selects coordinates inside it. Since dropping the ranking ($-0.0202$) or randomising it at equal budget ($-0.0124$ accuracy, $+0.0193$ WorstDrop) both degrade the result, the operative ingredient is \emph{retained-sensitivity ranking}, with structural overlap a property it recovers rather than a filter it needs; the intersection can thus be computed for analysis but skipped at request time. Second, a limitation: on CIFAR-10 every arm lands within $0.0035$ of full, so the arms separate only on the harder benchmark.

\subsection{The Intended Pretrained-Frozen PEFT Regime}
\label{subsec:pretrained}
Table~\ref{tab:pretrained} holds the PEFT premise fixed by giving LoRA and PALL-Adapter the same frozen ImageNet-pretrained ResNet-18. CIFAR inputs are converted from dataset normalization to the ImageNet statistics expected by the backbone; TinyImageNet already uses those statistics. All rows now use three seeds. PALL-Adapter reaches $0.9841$ on CIFAR-10, $0.8277$ on CIFAR-100, and $0.9047$ on TinyImageNet. The corresponding LoRA values are $0.9848$, $0.8215$, and $0.8980$. Thus the CIFAR-10 difference is negligible relative to seed spread, while Adapter is modestly higher on CIFAR-100 and TinyImageNet. Au must still be read by distance to chance, and these endpoint comparisons do not attribute the difference to the soft mask.

\begin{table}[t]
\centering
\small
\setlength{\tabcolsep}{1.2pt}
\begin{tabular}{@{}llccc@{}}
\toprule
\textbf{Dataset} & \textbf{Method} & \textbf{$A_{\mathrm{final}}$} & \textbf{W.Drop} & \textbf{Au} \\ \midrule
\multirow{2}{*}{C10} & LoRA & \textbf{0.9848${\pm}$.0154} & 0.0000${\pm}$.0000 & 0.5203${\pm}$.0966 \\
 & Adapter & 0.9841${\pm}$.0172 & \textbf{$-$0.0002${\pm}$.0003} & 0.5190${\pm}$.0988 \\ \midrule
\multirow{2}{*}{C100} & LoRA & 0.8215${\pm}$.0083 & 0.0213${\pm}$.0219 & 0.1613${\pm}$.0613 \\
 & Adapter & \textbf{0.8277${\pm}$.0061} & \textbf{0.0180${\pm}$.0020} & 0.1907${\pm}$.1066 \\ \midrule
\multirow{2}{*}{Tiny-IN} & LoRA & 0.8980${\pm}$.0049 & 0.0153${\pm}$.0110 & 0.0940${\pm}$.0092 \\
 & Adapter & \textbf{0.9047${\pm}$.0050} & \textbf{0.0127${\pm}$.0042} & 0.0940${\pm}$.0401 \\ \bottomrule
\end{tabular}
\caption{Corrected pretrained frozen-backbone regime (three seeds throughout). C10/C100/Tiny-IN denote CIFAR-10/CIFAR-100/TinyImageNet. CIFAR uses explicit ImageNet input normalization; TinyImageNet is already ImageNet-normalized. Chance Au is $0.5/0.2/0.1$. Bold marks the better endpoint for $A_{\mathrm{final}}$ and signed WorstDrop only; Au is interpreted by distance to chance.}
\label{tab:pretrained}
\end{table}

\subsection{PALL-Adapter Component Attribution}
\label{subsec:adapter-components}
Five matched arms isolate the request-time components under the corrected pretrained pipeline and schedules (three seeds): reset-only reinitializes the target adapter and classifier slice; reset+repair adds retained repair; uniform removes overlap protection from the shared update; mask/no-ascent removes the cached classifier ascent; full uses every component.

The RNG-neutral endpoints do not support the soft mask as the source of the adapter's behavior: full and uniform differ by at most $0.0006$ in $A_{\mathrm{final}}$ or WorstDrop, and full and mask/no-ascent are identical in both. The cached classifier ascent is the one component with a measurable, if small, effect: removing it leaves retention untouched but moves the forgotten task away from chance by $0.0040$ $[0.0020,0.0060]$ on CIFAR-100 and $0.0017$ $[0.0000,0.0030]$ on CIFAR-10, so it contributes to suppression without affecting retained accuracy. On CIFAR-100 reset-only is in fact better: retention $0.8324$ vs.\ $0.8244$, WorstDrop $0$ vs.\ $0.0147$, and a smaller mean distance to chance. Paired bootstrap intervals over three schedules include zero for the full--reset accuracy and Au-distance differences, whereas the WorstDrop difference favors reset-only. We therefore attribute the observed task suppression primarily to target-path reset; the overlap-aware shared update remains a theoretically delimited but empirically unsupported component at the measured endpoint.

\subsection{PALL Variants under the Overlap-Heavy Protocol}
\label{subsec:pall-overlap-heavy}
This separate from-scratch protocol uses three epochs per task and compares only the three PALL variants; it does not share the standard protocol's schedule or regime.

\begin{table}[t]
\centering
\footnotesize
\setlength{\tabcolsep}{2.5pt}
\begin{tabular}{@{}lccc@{}}
\toprule
\textbf{Method} & \textbf{$A_{\mathrm{final}}$} & \textbf{W.Drop} & \textbf{Au} \\ \midrule
\multicolumn{4}{@{}l}{\textit{CIFAR-10}} \\
PALL-Original & .9087${\pm}$.0266 & .0040${\pm}$.0069 & .5000${\pm}$.0000 \\
\textbf{EPALL} & .9152${\pm}$.0291 & .0008${\pm}$.0033 & .5000${\pm}$.0000 \\
\textbf{PALL-Adapter} & .7829${\pm}$.0610 & .0107${\pm}$.0101 & .4802${\pm}$.0816 \\ \addlinespace[2pt]
\multicolumn{4}{@{}l}{\textit{CIFAR-100}} \\
PALL-Original & .5013${\pm}$.0107 & .0187${\pm}$.0095 & .2033${\pm}$.0058 \\
\textbf{EPALL} & .5093${\pm}$.0148 & .0047${\pm}$.0050 & .2033${\pm}$.0058 \\
\textbf{PALL-Adapter} & .3284${\pm}$.0257 & .0427${\pm}$.0333 & .2147${\pm}$.0261 \\ \bottomrule
\end{tabular}
% Provenance: rows use the strict cifar10_main/cifar100_main tags.
\caption{PALL variants in the separate overlap-heavy from-scratch CIFAR regime (3 epochs/task; three seeds throughout). This protocol is not mixed with the standard comparison.}
\label{tab:pall_overlap_heavy}
\end{table}

As shown in Table~\ref{tab:pall_overlap_heavy}, \textbf{EPALL} has the highest retention accuracy among the three PALL variants on both CIFAR datasets ($0.9152$ vs.\ $0.9087$ on CIFAR-10 and $0.5093$ vs.\ $0.5013$ on CIFAR-100) while lowering WorstDrop from $0.0040$ to $0.0008$ and from $0.0187$ to $0.0047$ relative to PALL-Original. \textbf{PALL-Adapter} trades in the other direction: a smaller repair scope and Au close to chance, paid for with retention accuracy when training starts from scratch. Across-seed intervals overlap, so these three-seed results are descriptive trends rather than confirmed rankings.

\subsection{Nominal Retained-Repair Scope}
\label{subsec:efficiency}
PALL-Adapter reduces the logged nominal retained-repair scope relative to the full-network variants: from $0.0240$ to $0.0030$ on CIFAR-10 (an $87.5\%$ reduction), from $0.0210$ to $0.0055$ on CIFAR-100 (${\sim}74\%$), and from $0.0960$ to $0.0292$ on TinyImageNet (${\sim}70\%$). These ratios cover only the trainable repair scope: they do not establish lower end-to-end compute, peak memory, wall-clock time, or total parameter changes, and task-specific adapter storage still grows with the number of tasks. We treat them as implementation diagnostics and account resident storage separately next.

\subsection{Resident-State Accounting}
\label{subsec:storage}
Because the repair-scope ratio is not a storage measurement, Table~\ref{tab:storage} accounts resident tensor state directly on matched seed-0 schedules: EPALL keeps every task in one shared model with sparse masks and backups, whereas CLPU holds a full network per active task. CLPU consequently uses $1.6$--$2.2\times$ as much accounted state and grows $3.4$--$3.8\times$ faster per active task. The accounting covers resident parameters, side networks, masks, sparse backups, and replay/logits, and excludes optimizer state, activations, serialization, and interpreter overhead; it is an implementation-state comparison, not a disk-erasure or peak-memory claim.

\begin{table}[t]
\centering
\small
\setlength{\tabcolsep}{4pt}
\begin{tabular}{@{}llcc@{}}
\toprule
\textbf{Data} & \textbf{Method} & \textbf{Total (MiB)} & \textbf{MiB/task} \\ \midrule
\multirow{2}{*}{C10} & EPALL & \textbf{135.85} & \textbf{12.42} \\
 & CLPU & 212.94 & 42.59 \\ \midrule
\multirow{2}{*}{C100} & EPALL & \textbf{336.35} & \textbf{21.53} \\
 & CLPU & 730.78 & 81.20 \\ \bottomrule
\end{tabular}
\caption{Resident tensor-state accounting on matched seed-0 schedules (MiB, $2^{20}$ bytes) at maximum active tasks; MiB/task is the training growth per active task. Bold marks the smaller footprint. Excludes optimizer state, activations, serialization, and interpreter overhead.}
\label{tab:storage}
\end{table}

\subsection{Overlap-Response Benchmark}
\label{subsec:overlap-response}
The claim motivating EPALL is that protection matters \emph{as overlap grows}. We test it with a controlled benchmark: five pre-specified overlap levels $\times$ six methods $\times$ two datasets $=120$ matched runs, with the level set by schedule construction rather than measured post hoc (Figure~\ref{fig:overlap-response}).

\begin{figure}[t]
\centering
\includegraphics[width=.95\columnwidth]{Figures/aaai_overlap_response.pdf}
\caption{WorstDrop across five controlled overlap levels (vl to vh; mean over the matched runs at each level, lower is better). PALL-Original rises with overlap while EPALL stays comparatively flat; CLPU is identically zero because per-task networks cannot interfere. LoRA and the adapter are omitted: on a frozen random backbone their level reflects a capacity floor, not an overlap response.}
\label{fig:overlap-response}
\end{figure}

The benchmark separates the methods as the motivation predicts. On CIFAR-100, PALL-Original degrades monotonically with overlap, from $0.0130$ to $0.0390$ (a $3\times$ increase), whereas EPALL stays within $0.0060$--$0.0190$ throughout: the protection absorbs most of the extra interference. SalUn degrades similarly to PALL-Original, and CLPU is identically zero since per-task networks cannot interfere. CIFAR-10 preserves the ordering over a smaller range. Two caveats keep this descriptive: a level's realized overlap differs across methods, and pooled regressions on \emph{measured} overlap have wide intervals that mostly include zero (only the adapter's CIFAR-100 slope excludes it, $0.0398$ $[0.0190,0.0606]$). We report the level-wise trend, not a fitted dose--response law.

\paragraph{Overlap composition and persistence.} In the adapter runs with overlap logs the measured critical ratio $|S_{share\_crit}|/|S_{forget}|$ spans $0.35$--$0.56$ and the hard-protected fraction $0.07$--$0.11$. A read-only audit that follows each deleted task through the rest of the request history (two seeds; six trajectories per method and dataset) gives worst observed rebounds of $0.0490$/$0.0160$ for PALL-Adapter and $0.0185$/$0.0480$ for EPALL on CIFAR-10/100: suppression generally persists across later deletions but is not invariant. The schedules add no training request after deletion, so they do not test relapse under renewed learning as in UnCLe or CATA.

\paragraph{Anchor ablation.}
The $\ell_2$ anchor of Eq.~\eqref{eq:repair} defaults to the pre-forget weights $w_{old}$, which can themselves carry forgotten-task information. In the matched two-seed anchor ablation, replacing it with $w_{reinit}$ changes $A_{\mathrm{final}}$ from $0.9291$ to $0.9316$ and WorstDrop from $0.0045$ to $0.0037$ on CIFAR-10; on CIFAR-100 the corresponding changes are $0.4956$ to $0.5029$ and $0.0200$ to $0.0100$. This is directional retention evidence, not a privacy guarantee.

\paragraph{Sensitivity and importance criterion.}
The earlier pretrained $\{\alpha_f,\alpha_p,N_{iters}\}$ sweep used the legacy CIFAR-normalization pipeline and is excluded from the corrected endpoint evidence. In a two-seed CIFAR-10 sweep over bottleneck dimensions $\{4,8,16,32,64,128\}$, the endpoints were non-monotone; we retain $16$ as a compact fixed default rather than claim it is optimal. Replacing the gradient-magnitude criterion with the gradient-conflict energy $\mathrm{relu}(-g_{forget}\cdot g_{retain})$ helps the adapter on CIFAR-10 (WorstDrop $0.0160\!\to\!0.0063$) but degrades EPALL on CIFAR-100 ($0.0060\!\to\!0.0230$), so gradient magnitude remains the default.

\paragraph{Cross-architecture transfer.} On a ViT-T/8 backbone EPALL lowers mean WorstDrop on Split-CIFAR-100 from $0.0190$ to $0.0100$ while slightly improving accuracy ($0.3990$ vs.\ $0.3921$) and keeping Au near chance; on Split-CIFAR-10 the means are nearly equal. Transfer is dataset-dependent, with no significance claim.

\section{Limitations}
The request-driven interface requires the target task ID: it does not detect which task or records a request refers to, and does not cover class-incremental inference. Three-seed component controls find no endpoint benefit from PALL-Adapter's soft mask or classifier ascent beyond target-path reset, so that mechanism is exploratory rather than a second validated improvement. The direct EPALL controls separate the mechanism from its ablations only on CIFAR-100, and with three matched seeds the exact paired tests bottom out at $p=0.125$, so we report direction consistency and effect sizes instead of significance. A corrected, augmentation-matched membership-inference attack \cite{shokri2017membership,yeom2018privacy} with average-rank ties moves AUC from $0.5006{\pm}0.0205$ to $0.4959{\pm}0.0155$ on CIFAR-10 and from $0.5015{\pm}0.0169$ to $0.4998{\pm}0.0033$ on CIFAR-100, i.e.\ at chance before and after; but a loss-threshold attack is low-power by construction, and stronger calibrated attacks \cite{carlini2022membership} or formal auditing \cite{steinke2023privacy} would be needed to read this as removal. Against a same-method model never trained on the deleted task, EPALL's prediction agreement is $0.60$--$0.96$ on retained tasks but swings from $0.00$ to $0.98$ on the forgotten task across the four matched runs---again approximate, not exact, removal. Only the in-process rehearsal entry is deleted, not raw data, checkpoints, or external copies, and a frozen backbone can retain target information. Runtime is not hardware-normalized and adapter tuning is not automatic.

\section{Conclusion}
EPALL protects retention-sensitive shared weights, comes within half a point of CLPU on both standard Split-CIFAR benchmarks, and holds its advantage as controlled overlap grows. Component controls attribute PALL-Adapter's endpoint to target-path reset rather than its soft mask, so the adapter stays an audited exploratory variant. CLPU remains preferable when a full network per task is affordable, though its resident state is $1.6$--$2.2\times$ EPALL's. No result establishes exact, certified, or storage-level removal.


{\small
\bibliography{references}
}

\end{document}
